\section{常见实验问题与适用范围}
\label{sec:failure-boundaries}

只有把失败或异常结果与组成、实验条件及产物表征一起记录，才能判断方法的适用范围。本节讨论相竞争、污染与挥发、气氛与晶型、资料缺失及EHS限制。内容限于已发表条件，不把这些边界改写成新的实验方案。

\paragraph{相场与竞争。}
Ba--Y同类硅酸盐的系统助熔比较表明，成相结果必须结合组成和路线变量阅读，不能抽出一个热处理条件作为“万能条件”\cite{wierzbickawieczorek2017high}。后续研究通过非整比结构重审、陶瓷组成系列和多种衍射方法，揭示了调制结构、次生相以及竞争相或玻璃区域\cite{yamane2024synthesis,yamane2024microstructure,gulay2024navigation}。这些结果只界定各自考察的组成范围，不能证明未覆盖区域不存在其他相。机械化学研究仅提示预活化可能影响后续成相，不能把某个负结果解释为目标相的固有相界\cite{tzvetkov2001effects}。因此，表~\ref{tab:synthesis-conditions}中的产物、表征和合成条件必须来自同一文献，杂相或玻璃也不能脱离组成和热史单独解释。

\paragraph{污染、挥发与助熔边界。}
无ZnBa--Y陶瓷研究把次生相与玛瑙来源的\(\mathrm{SiO_2}\)污染联系到具体处理流程，说明表观配比偏离应先检查原料和处理过程，不能自动视为真实固溶证据\cite{yamane2024microstructure}。Y--Si--O的Czochralski研究可提供熔体、坩埚、气氛和提拉条件的参照，但现有资料不足以判断目标相是否发生Zn或Si挥发\cite{shoudu1999czochralski}。含铅、含氟助熔体系曾用于Zn--Si--O晶体生长，但这些历史条件只用于比较慢冷、分离和助熔剂作用，不能作为目标相的材料选择\cite{giess1982zn2sio4}。因此，应分别记录研磨介质带入、开放体系物质损失以及助熔剂残留或相选择；文献未给出烧前后组成或残留分析时，本文不指定具体失败机理。

\paragraph{气氛与晶型。}
BaZn$_2$Si$_2$O$_7$的低温单斜相与高温正交相由X射线和中子衍射区分，说明原位高温相与冷却后产物不可互换\cite{lin1999phase}。高温衍射还显示，Zn/Co取代以及Ba/Sr--Zn--Si/Ge固溶组成与相稳定或晶型边界联动；组成、温度和观察时点缺一项，都不足以定义稳定范围\cite{kerstan2013bazn2si2o7,thieme2016negative}。相关Ba--Zn--Si陶瓷的抗还原报道只说明该化合物在已发表条件下接受过气氛耐受考察，不能外推目标相的还原稳定性\cite{zou2019anti}；其他热膨胀与高温衍射比较也只作热致变化参照\cite{kerstan2012thermal}。

\paragraph{资料缺失与EHS。}
部分文献只给出合成方法或结构结论，没有同时报告配比、温度--时间、气氛、坩埚和冷却等完整信息\cite{ababaikeri2024ba5y12zn,ababaikeri2024ba3zn4si4o15,kaiser2002crystal,cai2024optimized,zhao2026crystal}。这表示本文尚未取得相应资料，不等于原文没有报告，也不等于实验没有该步骤；表~\ref{tab:synthesis-conditions}以\texttt{NA}保留这一边界。资料越不完整，可借鉴的具体条件越少。EHS边界同样只按已发表条件描述：含Ba粉体、高温挥发性物质、历史Pb/F助熔条件以及坩埚相容性和废物流均需单独审核，其中Pb/F路线尤其不能由历史实例改写为推荐\cite{giess1982zn2sio4}。任何复用均须经过独立的EHS、容器相容性和废物分流审核；本文不补写具体操作参数，也不对未报道的通风、防护、设备和废物分类作合规推定。

这四类边界不能合并成一个未经验证的失败机理。相竞争只在已发表的组成和热史范围内成立；污染或物质损失需要介质、烧前后组成或残留分析支持；气氛和晶型判断必须保留观察温度及冷却后的状态；资料缺失只说明当前认识的限度。因此，文献只报告“伴随”或只给出结构结果时，正文使用“伴随”或“一致”，不用“导致”。这样既保留负结果和成功记录，也避免把未报告项目、工艺异常和真实相区混为一谈。
